% RecSys Odyssey repository-backed lecture note.
% Add figures with: \coursefigure{figures/name.svg}{Caption}
\section{Why sequence order matters}

\begin{abstract}
The same events in a different order can describe a different intent.
\end{abstract}

\subsection{Concept}
A user who watches a trailer after browsing family films may want a film tonight; the reverse order may describe idle exploration. Sequential recommenders preserve event order and model transitions, recency and session boundaries. They complement rather than replace long-term preference: a useful state combines stable taste with the current task. The learning target is often the next item or next meaningful action.

\begin{equation*}
stateₜ = f(long-term profile, i₁, i₂, …, iₜ, contextₜ)
\end{equation*}

\subsection{Key terms}
\begin{itemize}
\item \textbf{Session.} A contiguous sequence of interactions representing one short-term intent.
\item \textbf{Next-item prediction.} Predicting the item likely to follow the observed prefix.
\item \textbf{Recency.} The principle that recent events often carry more current intent.
\end{itemize}
